\documentclass[12pt, a4paper]{article}
\usepackage[spanish]{babel}
\usepackage{color,graphicx}
\usepackage{graphics}
\usepackage{amsmath,amssymb,amsfonts,amsthm}
\usepackage[T1]{fontenc}
\usepackage{float}
\usepackage{xcolor}
\usepackage{amsthm}
\usepackage{tikz}
\usepackage{enumerate}
\usepackage{picinpar}
\usepackage{newtxtext,fourier}
\makeatother
\usepackage{thmtools}
\usepackage[framemethod=TikZ]{mdframed}
\mdfsetup{skipabove=1em,skipbelow=1em}

\newcommand{\sqed}{\hfill\includegraphics[height=2ex]{sancoya.png}}

\theoremstyle{definition}
\declaretheoremstyle[
    headfont=\bfseries\color{black!70!black}, bodyfont=\normalfont,
    mdframed={
        linewidth=0.7pt,
        rightline=true, topline=true, bottomline=true,
        linecolor=black, backgroundcolor=black!2!white,
    }
]{thmbox}

\declaretheorem[style=thmbox, name=Definición]{definition}
\declaretheorem[style=thmbox, numbered=no, name=Ejemplo]{example}
\declaretheorem[style=thmbox, name=Proposición]{prop}
\declaretheorem[style=thmbox, name=Teorema, numbered=no]{theorem}
\declaretheorem[style=thmbox, name=Lema]{lemma}
\declaretheorem[style=thmbox, numbered=no, name=Corolario]{corollary}
\declaretheorem[style=thmbox, numbered=no, name=Nota]{note}




\usepackage{geometry}
 \geometry{
 a4paper,
 total={170mm,260mm},
 left=25mm,
 top=30mm,
 }
\begin{document}
\setlength{\parindent}{0cm}
\hoffset-0.46cm
\voffset-1.46cm
\noindent
\hbox to \linewidth{
  \begin{minipage}[t]{40pt}
    \vspace*{0pt}
\includegraphics[scale=0.22]{Graphics/logo.eps}
  \end{minipage}
  \hss
  \begin{minipage}[t]{0.8\linewidth}
    \vspace*{15pt}\centering
     Universidad Nacional de Colombia\\
    \small Facultad de Ciencias\\
    Teoría analítica de números\\
    Parcial I
  \end{minipage}
  \hss
}
\phantom{xd}\\
\small
\textsf{Mateo Andrés Manosalva Amaris}
\normalsize
\dotfill
\vspace{0.3cm}

\textbf{Ejercicio 1.} Sea $f(x)$ una función definida para todo racional $x\in[0,1]$ y defina
\[
F(n)=\sum_{k=1}^{n} f\!\left(\frac{k}{n}\right),
\qquad
F_1(n)=\sum_{\substack{1\leq k\leq n\\ (k,n)=1}} f\!\left(\frac{k}{n}\right).
\]
\begin{enumerate}[(i)]
    \item Muestre que $F_1=\mu * F$ (convolución de Dirichlet).

    \begin{proof}
  Considere la suma con condición de coprimalidad, en efecto

  \begin{align*}
    F_1(n)&=\sum_{\substack{k=1\\(k,n)}}^n f\left(\frac{k}{n}\right)=\sum_{k=1}^n e((k,n))f\left(\frac{k}{n}\right)\\
    &=\sum_{k\leq n}\left(\sum_{\substack{d\mid k\\d\mid n}}\mu(d)\right)f\left(\frac{k}{n}\right)\\
    &=\sum_{d\mid n}\mu(d)\left(\sum_{\substack{k\leq n\\d\mid k}}f\left(\frac{k}{n}\right)\right) &&\text{Cambiando el orden de sumación}\\
    &=\sum_{d\mid n}\mu(d)\left(\sum_{dq\leq n}f\left(\frac{dq}{n}\right)\right)\\
    &=\sum_{d\mid n}\mu(d)\left(\sum_{q\leq n/d}f\left(\frac{q}{n/d}\right)\right)=\sum_{d\mid n}\mu(d)F\left(\frac{n}{d}\right)
  .\end{align*}

  Esto es, sencillamente que $F_1=\mu*F$, lo que se quería probar.
\end{proof}

    \item Pruebe que $\mu(n)$ es la suma de las raíces $n$-ésimas primitivas de la unidad:
    \[
    \mu(n)=\sum_{\substack{1\leq k\leq n\\(k,n)=1}} e^{2\pi i k/n}.
    \]

 \begin{proof}
      Sea $F(n)=\displaystyle\sum_{k=1}^ne^{2\pi ik/n}$, recordemos que

      $$F(n)=
\sum_{k=1}^n e^{2\pi i k/n}
= \left(\sum_{k=0}^{n-1} e^{2\pi i k/n}\right) + e^{2\pi i}
- 1
= \frac{1 - e^{2\pi i n/n}}{1 - e^{2\pi i/n}} + 1 - 1
= \frac{1 - e^{2\pi i}}{1 - e^{2\pi i/n}}
= 0 \quad (n>1),
    $$
es claro que $F(n)=1$ si $n=1$. Obtenemos que $F(n)=e(n)=(1*\mu)(n)$, por el numeral anterior $F_1=\mu* F=\mu*e=\mu$, esto es 

     \[F_1(n)=\sum_{\substack{1\leq k\leq n\\(k,n)=1}} e^{2\pi i k/n}=\mu(n).
    \]
    \end{proof}   
\end{enumerate}
\textbf{Ejercicio 2.} Sea
\[
\sigma_\alpha(n)=\sum_{d\mid n} d^\alpha.
\]
Pruebe:
\begin{enumerate}[(i)]
    \item Si $\alpha>0$ y $\alpha\neq 1$, entonces
    \[
    \sum_{n\leq x}\sigma_\alpha(n)
    =
    \zeta(\alpha+1)\frac{x^{\alpha+1}}{\alpha+1}
    +O\!\left(x^\beta\right),
    \qquad
    \beta=\max\{1,\alpha\}.
    \]

    \item Si $\alpha<0$ y $\beta=\max\{0,1+\alpha\}$, entonces para $x\geq 2$,
    \[
    \sum_{n\leq x}\sigma_\alpha(n)
    =
    \begin{cases}
    \zeta(1-\alpha)x+O(x^\beta), & \text{si } \alpha\neq -1,\\[1mm]
    \dfrac{\pi^2}{6}x+O(\log x), & \text{si } \alpha=-1.
    \end{cases}
    \]
\end{enumerate}

\begin{lemma}
    Sea $x\geq 1$, entonces tenemos las siguientes identidades

    \begin{enumerate}[$\bullet$]
        \item 
        $\displaystyle\sum_{n \le x} \frac{1}{n^s}= \frac{x^{1-s}}{1-s} + \zeta(s) + O(x^{-s})\quad \text{si } s > 0,\; s \ne 1. $
        \item $\displaystyle\sum_{n \le x} n^\alpha
= \frac{x^{\alpha+1}}{\alpha+1} + O(x^\alpha)
\quad \text{si } \alpha \ge 0.$
    \end{enumerate}
\end{lemma}

\begin{proof}
Note que
\begin{align*}
\sum_{n \le x} \frac{1}{n^s}
&= 1 + \int_1^x \frac{1}{t^s} \, d[t] \\
&= 1 + \int_1^x \frac{1}{t^s} \, dt - \int_1^x \frac{1}{t^s} \, d\{t\} \\
&= 1 + \frac{x^{1-s}}{1-s} - \frac{1}{1-s} - \int_1^x \frac{1}{t^s} \, d\{t\} \\
&= 1 + \frac{x^{1-s}}{1-s} - \frac{1}{1-s}
- \left( \frac{\{x\}}{x^s} + s \int_1^x t^{-s-1} \{t\} \, dt \right).
\end{align*}

Como $\{t\} = O(1)$, entonces
\begin{align*}
\sum_{n \le x} \frac{1}{n^s}
&= \frac{x^{1-s}}{1-s} + 1 + \frac{1}{s-1}
- s \int_1^x \frac{\{t\}}{t^{s+1}} \, dt + O(x^{-s}) \\
&= \frac{x^{1-s}}{1-s} + \frac{s}{s-1}
- s \int_1^x \frac{\{t\}}{t^{s+1}} \, dt + O(x^{-s}) \\
&= \frac{x^{1-s}}{1-s} + \zeta(s) + O(x^{-s}),
\end{align*}

por la representación integral de $\zeta(s)$, esta representación integral tiene sentido para $s > 0$
con $s \neq 1$, lo que concluye el resultado.\\

Para la segunda identidad, procedemos de la misma forma

\begin{align*}
\sum_{n \le x} n^{\alpha}
&= 1 + \int_1^x t^\alpha \, d[t] \\
&= [x]x^\alpha - \int_1^x [t] \, d(t^\alpha) \\
&= [x]x^\alpha - \alpha \int_1^x [t] t^{\alpha-1} \, dt \\
&= x^{\alpha+1} + O(x^\alpha)
   - \alpha \int_1^x (t + O(1)) t^{\alpha-1} \, dt \\
&= x^{\alpha+1} + O(x^\alpha)
   - \alpha \int_1^x t^\alpha \, dt
   + O\!\left(\int_1^x t^{\alpha-1} dt\right) \\
&= \frac{\alpha x^{\alpha+1}+x^{\alpha+1} }{\alpha+1}+ O(x^\alpha)- \alpha \frac{x^{\alpha+1}-1}{\alpha+1}+ O(x^\alpha) \\
&= \frac{x^{\alpha+1}}{\alpha+1} + O(x^\alpha).
\end{align*}
\end{proof}

Con estas estimaciones podemos proceder con la prueba de manera más directa.\\

\begin{proof}
(i) Por el lema anterior

\begin{align*}
\sum_{n \le x} \sigma_\alpha(n)
&= \sum_{n \le x} \sum_{m \mid n} m^\alpha
= \sum_{d \le x} \sum_{m \le x/d} m^\alpha \\
&= \sum_{d \le x} \left(
\frac{1}{\alpha+1}\left(\frac{x}{d}\right)^{\alpha+1}
+ O\!\left(\frac{x^\alpha}{d^\alpha}\right)
\right) \\
&= \frac{x^{\alpha+1}}{\alpha+1}
\sum_{d \le x} \frac{1}{d^{\alpha+1}}
+ O\!\left(x^\alpha \sum_{d \le x} \frac{1}{d^\alpha}\right) \\
&= \frac{x^{\alpha+1}}{\alpha+1}
\left(
\frac{x^{-\alpha}}{-\alpha} + \zeta(\alpha+1) + O\left(x^{-\alpha-1}\right)
\right) + O\left(
x^\alpha \left(
\frac{x^{1-\alpha}}{1-\alpha} + \zeta(\alpha) + O\left(x^{-\alpha}\right)
\right)
\right) \\
&= \frac{\zeta(\alpha+1)}{\alpha+1} x^{\alpha+1}
+ O(x) + O(1) + O(x^\alpha) \\
&= \frac{\zeta(\alpha+1)}{\alpha+1} x^{\alpha+1}
+ O(x^\beta),
\end{align*}

donde $\beta = \max\{1,\alpha\}.$\\

(ii) De la misma forma

\[
\sum_{n \le x} \sigma_{\alpha}(n)
= \sum_{n \le x} \sum_{d \mid n} d^{\alpha}
= \sum_{d \le x} d^{\alpha} \sum_{\substack{n \le x \\ d \mid n}} 1
= \sum_{d \le x} d^{\alpha} \left( \frac{x}{d} + O(1) \right)= x \sum_{d \le x} d^{\alpha-1} + O\!\left(\sum_{d \le x} d^{\alpha}\right).
\]

 Si $\alpha \neq -1$, entonces

\begin{align*}
    \sum_{n \le x} \sigma_{\alpha}(n)&=x\left(\frac{x^{\alpha}}{\alpha}+\zeta(1-\alpha)+O(x^{\alpha-1})\right)+O\left(\frac{x^{1+\alpha}}{1+\alpha}+\zeta(-\alpha)+O(x^\alpha)\right)\\
    &=x\zeta(1-\alpha)+O(x^{1+\alpha})+O(1)
.\end{align*}

Tenemos dos consideraciones, en efecto si $-1<\alpha<0$, entonces $O(x^{1+\alpha})$ domina a $O(1)$, si $\alpha<-1$, domina $O(1)$. así pues concluimos 

$$\sum_{n \le x} \sigma_{\alpha}(n)=x\zeta(1-\alpha)+O(x^{\beta})\quad \text{donde } \beta=\max\{0,1+\alpha\}.$$

Para concluir, en el caso $\alpha=-1$,
\begin{align*}
\sum_{n\le x}\sigma_{-1}(n)
&=\sum_{n\le x}\sum_{d\mid n}\frac{1}{d}
=\sum_{d\le x}\frac{1}{d}\sum_{\substack{n\le x\\ d\mid n}}1
=\sum_{d\le x}\frac{1}{d}\left\lfloor \frac{x}{d}\right\rfloor \\
&=\sum_{d\le x}\frac{1}{d}\left(\frac{x}{d}+O(1)\right)\\
&= x\sum_{d\le x}\frac{1}{d^2} + O\!\left(\sum_{d\le x}\frac{1}{d}\right) \\
&= x\left(\frac{\pi^2}{6}+O\!\left(\frac{1}{x}\right)\right) + O(\log x)
= \frac{\pi^2}{6}\,x + O(\log x).
\end{align*}
\end{proof}

\textbf{Ejercicio 3:} Sea

$$F(s)=\sum_{n=1}^\infty \dfrac{a_n}{n^s},$$

donde $a_n$ es completamente multiplicativa y la serie converge absolutamente para $\sigma>\sigma_a$. Probar que para $\sigma>\sigma_a$ se cumple que 

$$\dfrac{F^\prime(s)}{F(s)}=-\sum_{n=1}^\infty\dfrac{a_n\Lambda(n)}{n^s}$$

\begin{lemma}
    Sean $f$ y $g$ las funciones aritmética asociadas a $F(s)$ y $G(s)$, considere $h=f*g$ y $H(s)$ su serie de Dirichlet asociada. Si $F(s)$ y $G(s)$ convergen absolutamente en $s$, entonces $H(s)$ también y además $H(s)=F(s)G(s)$ 
\end{lemma}

\begin{proof}
     Por la convergencia absoluta
     \begin{align*}
         F(s)G(s)&=\sum_{k=1}^\infty\dfrac{f(k)}{k^s}\sum_{m=1}^\infty\dfrac{g(m)}{m^s}\\
         &=\sum_{k=1}^\infty\sum_{m=1}^\infty\dfrac{f(k)g(m)}{k^sm^s}\\
         &=\sum_{n=1}^{\infty}\dfrac{1}{n^s}\sum_{km=n}f(k)g(m)\\
         &=\sum_{n=1}^\infty\dfrac{(f*g)(n)}{n^s}
     \end{align*}

     La convergencia absoluta de $H(s)$ se sigue de 

     $$\sum_{n=1}^\infty\left|\dfrac{h(n)}{n^s}\right|=\sum_{n=1}^{\infty}\dfrac{1}{|n^s|}\sum_{km=n}|f(k)||g(m)|=\sum_{k=1}^\infty\left|\dfrac{f(k)}{k^s}\right|\sum_{m=1}^\infty\left|\dfrac{g(m)}{m^s}\right|<\infty$$
\end{proof}

\begin{lemma}
    Sea $f$ una función completamente multiplicativa, entonces $(f(g*h))(n)=(fg*fh)(n)$.
\end{lemma}

\begin{proof}
Por hipótesis, $f$ es completamente multiplicativa, esto es, $f(n)=f(j)f\left(\dfrac{n}{j}\right)$, así pues
    \begin{align*}
        (f(g*h))(n)&=f(n)\sum_{j\mid n}g(j)h\left(\dfrac{n}{j}\right)\\
        &=\sum_{j\mid n}g(j)f(j)h\left(\dfrac{n}{j}\right)f\left(\dfrac{n}{j}\right)\\
        &=\sum_{j\mid n}fg(j)fh\left(\dfrac{n}{j}\right)\\
        &=(fg*fh)(n).
    \end{align*}
\end{proof}

Ahora procedemos con la prueba.

 \begin{proof}
     En efecto, para $\sigma>\sigma_a$ se cumple que

    \begin{align*}
        F^\prime(s)&=-\sum_{n=1}^\infty\dfrac{a_n\log n}{n^s}\\
        &=-\sum_{n=1}^\infty\dfrac{a_n(\Lambda*1)(n)}{n^s} &&\text{Teorema Fundamental de la Aritmética}\\
        &=-\sum_{n=1}^\infty\dfrac{a_n\Lambda(n)*a_n}{n^s}\\
        &=-\sum_{n=1}^\infty\dfrac{a_n\Lambda(n)}{n^s}\sum_{n=1}^\infty\dfrac{a_n}{n^s}=-\sum_{n=1}^\infty\dfrac{a_n\Lambda(n)}{n^s}F(s),
    \end{align*}
    donde la separación de esta convolución como un producto, está justificada por el lema anterior, ya que $F(s)$ converge absolutamente en el semiplano $\sigma>\sigma_a$, además como $\Lambda(n)\leq \log n$ y $F^\prime(s)$ converge absolutamente en $\sigma>\sigma_a$, entonces también lo hace $\displaystyle\sum_{n=1}^\infty\dfrac{a_n\Lambda(n)}{n^s}.$ De esto se sigue el resultado.
 \end{proof}


\textbf{Ejercicio 4.} Sea
\[
F(s)=\sum_{n=1}^{\infty} \frac{f(n)}{n^{s}},
\]
convergente para $\sigma>c$, donde $c$ es finito, y suponga que $f(n)\ge 0$ para todo $n\ge n_0$.

\begin{enumerate}[(a)]
    \item Demostrar que si $F(s)$ es analítica en un disco alrededor de $s=c$, entonces la serie converge en $\sigma>c-\varepsilon$ para algún $\varepsilon>0$.

    \item Concluir que si la abscisa de convergencia $\sigma_c$ es finita, entonces $F(s)$ tiene una singularidad en el punto $s=\sigma_c$ sobre el eje real.
\end{enumerate}

\begin{proof}
    Sea $a=1+c$, como $F$ es analítica en $a$, entonces por el teorema de Taylor

    \begin{align}
        F(s)=\sum_{k=0}^\infty\dfrac{F^{(k)}(a)}{k!}(s-a)^k,
    \end{align}

 además el radio de convergencia de esta serie de potencias es mayor que 1, ya que $F$ es analítica en $c$. Como vimos $\displaystyle F^\prime(s)=-\sum_{n=1}^\infty\dfrac{f(n)\log n}{n^s}$, luego derivando de forma iterada, obtenemos que

    $$F^{(k)}(a)=(-1)^k \sum_{n=1}^{\infty}f(n)(\log n)^kn^{-a}$$

podemos reescribir (1) como:

\begin{align}
F(s)=\sum_{k=0}^{\infty}\sum_{n=1}^\infty\dfrac{(a-s)^k}{k!}f(n)(\log n)^k n^{-a}
\end{align}

Nuevamente, como el radio de convergencia de (2) es mayor que 1, entonces, es convergente para algún $s\in \mathbb{R}$ tal que $s=c-\varepsilon$, con $\varepsilon>0$. En efecto $a-s=1+\varepsilon$, así pues, la serie (2) tiene términos no negativos para $n\geq n_0$, por lo tanto, el teorema de Tonelli permite cambiar el orden de sumación. Obtenemos que

$$F(c-\varepsilon)=\sum_{n=1}^\infty\dfrac{f(n)}{n^a}\sum_{k=0}^\infty\dfrac{((1+\varepsilon)\log n)^k}{k!}=\sum_{n=1}^\infty\dfrac{f(n)}{n^a}e^{(1+\varepsilon)\log n}=\sum_{n=1}^\infty \dfrac{f(n)}{n^{c-\varepsilon}}.$$

Esto es, $F(s)=\displaystyle\sum_{n=1}^\infty f(n)n^{-s}$ converge para $s=c-\varepsilon$, entonces en particular, converge en el semiplano $\sigma>c-\varepsilon$.
\end{proof}

Finalmente, supongamos que $\sigma_c$ no es una singularidad de $F$, luego $F(s)$ es analítica en $s=\sigma_c$, pero por lo anterior, esto significa que $F(s)$ converge en $\sigma>c-\varepsilon$ lo que contradice que $\sigma_c$ es la abscisa de convergencia.

\end{document}